% One-page HOA coexistence note (course extension; not an official 2023 A page).
\documentclass[11pt]{letter}
\usepackage[margin=1in]{geometry}
\usepackage[hidelinks]{hyperref}

\signature{GoHiMCM 2023 A study team}
\address{Modeling desk\\GoHiMCM}

\begin{document}
\begin{letter}{Homeowners Association Board\\One-Hectare Common Lawn}

\opening{Dear Board Members,}

You asked what happens if one puffball-stage dandelion sits just outside the
west fence of the common hectare, and what a sane mowing rule looks like.
We simulated a year of wind, temperature, and soil moisture as a
\emph{temperate lawn scenario} (not a weather-station download). Occupied
lawn---cells with at least $0.05$ plants per square meter---is about
$0.21$ of the hectare by month 12 if nobody mows
(\texttt{results/monthly\_metrics.csv}). That is a patch growing in from the
west, not a uniform yellow carpet. Most of the head-count is seedlings;
flowering adults at year-end are a few thousand, not a million.

Water stress matters. Cutting the soil-moisture index in half (a drought
scenario) lowers month-12 occupancy to about $0.12$.
A hotter, drier climate in our arid table ends the year near $0.027$.
A tropical table with wind blowing \emph{away} from the lawn barely
invades ($0.006$ occupancy), because seeds leave the property.

If the goal is a tidy lawn rather than eradication of a useful edible plant,
you do not need weekly scalping. On our cost--occupancy curve
(\texttt{results/bioeconomic\_pareto.csv}), mowing every two weeks at full
height reduction reaches about $0.008$ occupancy at relative cost $39$.
Mowing every week costs twice as much for the same occupancy and is wasted
effort. Leaving the stand and harvesting greens is consistent with the
low impact score dandelion receives against kudzu and Japanese knotweed
in our three-species comparison: those two are the board's real ``do not
plant / do not ignore'' species if they appear on the edges of the
development.

We recommend: (1)~keep a two-week mowing cadence through the spring and fall
bloom windows; (2)~do not spray the whole hectare for one west-edge patch;
(3)~revisit the rule after a drought summer, when natural cover is already
lower.

\closing{Respectfully,}
\end{letter}
\end{document}
